\section{Inflation and Deflation}

\label{sec:deflation}
% ============================================================================

HANZO has a \textbf{fixed supply with deflationary dynamics}. No new tokens are
minted after genesis. The only supply-reducing mechanism is the fee burn.

\subsection{Burn Mechanism}

A fraction $\zeta = 0.25$ of all platform fees collected in HANZO is
permanently burned (sent to a provably unspendable address). Let $F(t)$ denote
cumulative fees at time $t$. The effective circulating supply is:
\begin{equation}
S_{\text{eff}}(t) = S_{\max} - \zeta \cdot F(t) - S_{\text{staked}}(t) -
S_{\text{locked}}(t),
\end{equation}
where $S_{\text{staked}}(t)$ is the total staked amount and $S_{\text{locked}}(t)$
is the total in fee-reduction contracts.

\subsection{Equilibrium Analysis}

\begin{definition}[Burn-Stake Equilibrium]
A state $(p^*, S_{\text{staked}}^*, S_{\text{locked}}^*)$ is a burn-stake
equilibrium if the token price $p^*$ is such that:
\begin{enumerate}[leftmargin=1.4em]
    \item Node operators find staking profitable:
    $\frac{r_i}{s_i} \geq r_{\text{alt}}$, where $r_{\text{alt}}$ is the
    alternative cost of capital.
    \item Compute consumers find holding for fee reduction profitable:
    the fee savings exceed the opportunity cost of locked capital.
    \item The burn rate does not exceed new demand: $\zeta \cdot F'(t) <
    D'(t) \cdot p^*$.
\end{enumerate}
\end{definition}

Under mild assumptions (monotonically growing compute demand with $\gamma > 0$
and bounded staking ratio $S_{\text{staked}}/S_{\max} < 0.6$), the system
converges to a stable equilibrium where increasing demand raises the token price,
which increases the dollar-denominated staking reward, attracting more stakers,
which reduces circulating supply, further supporting the price. The burn
mechanism provides a permanent supply floor that rises with cumulative usage.

\subsection{Projected Burn}

At observed fee rates (30 bps market fee on \$2.4M daily compute volume at
steady state), the annual burn is approximately:
\begin{equation}
B_{\text{annual}} = 0.25 \times 0.003 \times 2{,}400{,}000 \times 365 /
p \approx \frac{657{,}000}{p} \text{ HANZO},
\end{equation}
where $p$ is the token price in USD. At $p = \$0.10$, this burns 6.57M tokens
per year (0.66\% of total supply). At $p = \$1.00$, the burn is 657K tokens
(0.066\%). The deflationary effect is strongest at low prices, providing a
natural floor.

% ============================================================================
